{\small\sffamily \textbf{Abstract.} \articleABSTRACT\par}
\smallskip
{\small\sffamily A partial replication of Almeida~\cite{almeida2023}.\par}

\section{Introduction}

A Bloom filter gives a fast answer to a membership query. It can return a
false positive for an absent key. The probability usually reported is an
average over absent keys. If one key is tested against many filters, the
probability for that key matters instead. In a standard Bloom filter,
several hash probes can land on the same bit. A repeated probe
leaves one fewer bit for the query to check.

Almeida~\cite{almeida2023} compares standard and partitioned Bloom filters at finite sizes.
Under independent uniform hashing, four partitioned-to-standard mean
false-positive ratios in his Table~2 exceed an upper bound valid for every
common positive integer insertion count in both layouts. These are the
quarter- and half-nominal-load values at 64 bits with four or eight probes.
Integer rounding cannot explain these printed values. Almeida reports a loop
error and a matching corrected example (personal communication, 7 November
2026). We have not independently reproduced his diagnosis.

This partial replication covers four tables and two curves in Sections~3
and~4, plus three double-hashing mechanisms in Section~5. It recovers
nominal-load mean rates and fixed-key collision effects. It matches 99
of 124 original printed cell occurrences at the reported precision,
counting repeated ratios in each source table.
Table~\ref{tab:scope} summarizes the outcomes.

\begin{table}[tbp]
\centering
\caption{Scope and outcome. Source table and figure numbers refer to
Almeida~\cite{almeida2023}; section references refer to this replication.}
\label{tab:scope}
\begin{tabularx}{\linewidth}{@{}>{\raggedright\arraybackslash}p{0.22\linewidth}>{\raggedright\arraybackslash}X@{}}
\toprule
Status & Evidence and result \\
\midrule
Replicated conclusions & Nominal-load mean rates and collision effects
(source Tables~1, 3, and~4; Figures~3 and~4; Sections~\ref{sec:mean}
and~\ref{sec:fixed}). The three double-hashing mechanisms are recovered
(source Figures~5--7; Section~\ref{sec:double}). \\
Original-table differences & 25 differences in Tables~1--3, including four
Table~2 ratios above the bound at 64 bits. Almeida reports a loop error and
a matching example; we have not independently reproduced his diagnosis
(Section~\ref{sec:mean}). \\
Exact extension & Uniform hash-pair enumeration at 512 bits, eight probes,
and 44 insertions gives mean rates of 0.649\% (standard) and 1.442\%
(partitioned) (Section~\ref{sec:double}). Source Figures~5--7 motivate
this calculation but do not report these rates. \\
\bottomrule
\end{tabularx}
\end{table}

\section{Methods}

A standard filter has one array of $m$ bits. Each insertion sets $k$
positions and a query is positive if all its positions are set. In a
partitioned filter, $k$ divides $m$; each of $k$ equal parts has $s=m/k$
bits and receives one position per key. A collision between two positions
within the index range of a part leaves two different physical bits because the positions
belong to different parts.

\subsection{Independent hashing}

Under independent uniform hashing, we use three global false-positive rates
for $n$ insertions:
\begin{itemize}[noitemsep,topsep=2pt,leftmargin=*]
\item $F_a(n,m,k)$ is the approximate rate for the standard filter.
\item $F_s(n,m,k)$ is the exact rate for the standard filter.
\item $F_p(n,m,k)$ is the exact rate for the partitioned filter.
\end{itemize}
Each global rate averages over insertion hash outputs and an independent
absent query. Both layouts use the same nonnegative integer $n$.
Arguments are omitted when fixed.

Fix a query that requires $d$ distinct physical bits. After $t$ insertion
probes, let $p_t(j)$ be the probability that exactly $j$ of those bits
have been set. Start with $p_0(0)=1$. A new independent uniform probe
sets a previously unset query bit with probability $(d-j)/m$.
The probability of satisfying the query after $kn$ probes is
$p_{kn}(d)$. This gives the same quantity as Almeida's conditional
false-positive probability, using a recurrence on the number of covered
query bits. This fixed-key conditional rate averages over insertions
with the query's hash outputs held fixed, unlike the global rate $F_s$.

The program runs one chain with $d=k$. For a query that needs $d$ bits,
it averages $\binom{j}{d}/\binom{k}{d}$ over the chain's final state.
The ratio is the chance that the $d$ required bits are among the $j$
covered bits. To obtain the global rate, the program averages over
the number of distinct positions in $k$ independent query probes. The partitioned
rate under independent hashing is
$F_p=[1-(1-1/s)^n]^k$. Each part has one required position and its insertion
draws are independent of the other parts. The approximate standard
formula in the source raises expected bit occupancy to the $k$th power:
$F_a=[1-(1-1/m)^{kn}]^k$.

\subsection{Naive double hashing}

Two independent uniform residues $(a,b)$ per key generate the sequence
$a+ib$ modulo $m$, or modulo $s$ inside each partition. Pairs for distinct
keys are independent, and the absent query's pair is independent of the
insertion pairs. The probes generated by one pair are generally dependent,
so this joint model differs from the independent-hashing model defining
$F_s$ and $F_p$. A standard filter treats the sequence as an unordered
set of bits. A partitioned filter keeps the index $i$ as part of the
physical address.

To compute an exact rate for a fixed query, enumerate all possible
insertion pairs. For each subset $U$ of the query's required bits, count
how many insertion pairs avoid every bit in $U$. The chance that all
$n$ insertions avoid $U$ is the $n$th power of that fraction. Summing over
$U$ with sign $(-1)^{|U|}$ gives the chance that every required bit has
been set. The calculation retains dependence among the $k$ bits
within each insertion pair.

\subsection{Comparison protocol}

For $m$ bits, $k$ probes, and a fraction $f$ of nominal capacity, the
primary analysis uses $n=\lfloor f(m/k)\log 2\rfloor$ insertions. The
source gives the nominal capacity as a real number. The printed $F_a$ values in Table~1 identify the floor count in all
eight nominal-load configurations.
For reduced load, the floor is an explicit interpretation. A computed
cell matches if it rounds to the published value, allowing $10^{-12}$
for numerical error at the boundary.

An integer inclusion-exclusion calculation checks the recurrence across
the table grid. Tiny exhaustive enumerations check both layouts under
independent and double hashing. The direct simulation builds 120,000
independent filters for each load at $m=64,k=4$. Within each filter it
tests one fixed query for each distinct-bit count, from one to four. At
full load it also tests a fresh random query per standard filter and
runs a separate set of 120,000 partitioned filters. Wilson intervals
describe each marginal probability. The queries share filters within
a trial, so the intervals are neither independent nor simultaneous.

\section{Average false-positive rates}
\label{sec:mean}

At nominal capacity, the independent calculation matches all 24
$F_a$, $F_s$, and $F_p$ cells in Almeida's Table~1 at their printed
precision. Seven of the eight ratios match. The exception is
$m=64,k=8$, where the paper gives 1.21703762 and the exact calculation
gives 1.21703636. The reconstructed size sweep (Figure~\ref{fig:mean})
shows the same broad pattern as the source's Figure~3. The ratio
approaches one as the filter grows. With $k=16$, the ratio is about
1.098 at $m=512$ and about 1.011 at $m=4096$.

\input{generated/mean-figure.tex}

The reduced-load differences are larger. Table~\ref{tab:comparison}
shows the four largest Table~2 discrepancies, all at $m=64$. Seven of the
16 reduced-load values in the source table differ at the printed
precision. The eighth mismatch in the table is the nominal-load
ratio repeated from Table~1. Switching from floor to nearest or ceiling
reconciles none of the printed exceptions.
\input{generated/comparison-table.tex}

Under independent uniform hashing, no common positive integer insertion
count $n$ in both layouts produces these four ratios. Set
$a=(1-1/m)^k$ and $b=1-k/m$. Then $0\le b\le a<1$.
For every integer $n\ge1$, the lower bound $F_a\le F_s$ gives
\[
  \frac{F_p}{F_s}
  \le \frac{F_p}{F_a}
  =\left(\frac{1-b}{1-a}
     \frac{\sum_{j=0}^{n-1} b^j}{\sum_{j=0}^{n-1} a^j}\right)^k
  \le \left(\frac{1-b}{1-a}\right)^k.
\]
At $m=64$, the bound is 1.09842 for four probes and 1.54612 for eight,
independent of $n$. The reported half- and quarter-load ratios exceed
these bounds.

Almeida (personal communication, 7 November 2026) reports an off-by-one
error in a loop using $n$ that stopped one iteration early. For
$m=64,k=4,n=2$, his corrected standard false-positive rate is
$0.00020343338632548505$, with ratio $1.0569749279689804$.
His ratio agrees numerically with our independent value,
$1.0569749279689811$. No author code or corrected tables were supplied.
We have not inspected his code or independently reproduced the reported error.
The matching example does not establish that all 25 differences against the
original published tables are resolved.

\section{False positives for a fixed key}
\label{sec:fixed}

The conditional rate rises as repeated probes reduce the number of
distinct bits the query needs. Thirty-two of the 48 cells in Almeida's
Table~3 match, including all 16 at nominal load. All 20 cells of
Table~4 match. The mismatches in Table~3 occur at reduced load. For
example, at $m=64,k=4$ and quarter load, a one-bit query has a rate
581.88 times the global average in this calculation, compared with
804.25 in the paper. In the reconstructed model the one-bit rate is
11.84\%, and the global rate is about 0.0203\%.

The tail plot (Figure~\ref{fig:tail}) reconstructs the source's
Figure~4. For $m=512,k=16$ at nominal load, at least three collisions
occur for about one query in 914. For these queries, the false-positive
rate is at least 6.78 times the global rate. At least five collisions
occur for about one in 1.34 million, with a rate of about 28.30 times
the global rate. These thresholds refer to collision classes of absent
queries under independent uniform hashing.

\input{generated/tail-figure.tex}

\input{generated/simulation-results.tex}

\section{Double hashing}
\label{sec:double}

At $m=512,k=8$, a standard filter with step $b=0$, $256$, or $128$
touches one, two, or four physical bits, respectively. Using zero, half,
or a quarter of the part size as the step in a partitioned filter still
touches eight bits. The pairs $(7,5)$ and $(42,-5)$ visit the same
standard-filter bits in reverse order. Their partitioned addresses have no overlap. The
pairs $(7,5)$ and $(12,5)$ share seven standard-filter positions and
no partitioned positions. These three
examples recover the mechanisms in Almeida's Figures~5--7.

The enumeration measures how these cases affect false positives.
For a standard filter, translation removes the starting residue.
Multiplication by an odd unit modulo 512 preserves the insertion-pair
law and maps steps with the same greatest common divisor to each
other. Each step class thus has one fixed-key rate.
For a partitioned query with pair $(a_q,b_q)$, subtracting the query's
pair from each insertion pair maps the query to $(0,0)$ while
preserving the insertion law. Every fixed query therefore has the same rate
in this model.

\input{generated/double-hash-table.tex}

At $n=44$, the zero-step standard query has a 49.81\% false-positive
rate and accounts for one in 512 hash pairs. The average standard
double-hashing rate is 0.649\%. Every fixed partitioned query has a
1.442\% rate, which also gives its global rate. Partitioning removes the
step-dependent variation in this model and raises the average at this
configuration. The 0.390\% partitioned result under independent hashing
uses a different joint distribution of insertion probes.

\section{Scope and availability}

The artifact retains all 25 differences against the original published
tables, with each source value beside its computed value, model, and
integer count. We have not independently reproduced Almeida's reported
diagnosis (Section~\ref{sec:mean}).

The results concern uniform hash outputs at finite filter sizes. The
experiment does not measure deployed hash functions, workload-dependent
key frequencies, SIMD throughput, or the applications in Sections~6 and
7 of the source. The double-hashing extension covers one allocation and
one load. The comparison values come from the public accepted manuscript,
arXiv:2009.11789v2. Almeida responded on 7 November 2026. The calculation
uses no author code or experiments.

The companion artifact contains the code, transcribed values, and
generated tables and plots. A single Python command reruns the
experiment offline. The PDF build reads the generated results.
